\documentclass[11pt]{article}
\usepackage[margin=1in]{geometry}
\usepackage{amsmath,amssymb,amsthm,mathtools}
\usepackage{verbatim}
\usepackage{hyperref}
\usepackage{booktabs}

\newtheorem{theorem}{Theorem}
\newtheorem{lemma}[theorem]{Lemma}
\newtheorem{proposition}[theorem]{Proposition}
\newtheorem{corollary}[theorem]{Corollary}
\newtheorem{remark}[theorem]{Remark}
\newtheorem{conjecture}[theorem]{Conjecture}
\theoremstyle{definition}
\newtheorem{example}[theorem]{Example}
\newtheorem{definition}[theorem]{Definition}

\newcommand{\Fock}{\mathcal{F}}
\newcommand{\Uq}{U_q(\widehat{\mathfrak{sl}}_e)}
\newcommand{\Lam}{\Lambda_0}
\newcommand{\vke}{v_{k',e}}
\newcommand{\ZZ}{\mathbb{Z}}
\newcommand{\QQ}{\mathbb{Q}}
\newcommand{\supp}{\operatorname{supp}}
\newcommand{\Lcr}{\mathcal{L}}
\newcommand{\ee}{\tilde e}
\newcommand{\ff}{\tilde f}
\newcommand{\eps}{\varepsilon}
\newcommand{\Aa}{{\mathsf A}}
\newcommand{\Rr}{{\mathsf R}}
\newcommand{\Ce}{C_e^{(1)}}

\title{At level $1$, the Kashiwara-crystal string of $\vke = P_e^{k'}|\emptyset\rangle$
       has length at most one:\\ $\eps_i(\vke) \le 1$ for all $i \in \ZZ/e\ZZ$}
\author{Clio}
\date{2026-08-11}

\begin{document}
\maketitle

\begin{abstract}
Let $e \ge 2$ and let $\Fock_e$ denote the level-$1$ Uglov $q$-Fock space of
$\Uq$, with standard basis $\{|\lambda\rangle\}$ indexed by partitions.
Let $P_e$ be the ribbon-sign operator (Leclerc--Thibon), and let
$\vke := P_e^{k'}|\emptyset\rangle$. Let $(\Lcr, \{[\mu]\})$ denote the
Kashiwara crystal lattice of $\Fock_e$, and let $\{\ee_i, \ff_i\}_{i \in \ZZ/e\ZZ}$
be the Kashiwara operators.

We prove:
\begin{itemize}
  \item[(C5, $\ell=1$)] \emph{Kashiwara string-length bound.} For every $e \ge 2$,
  every $k' \ge 1$, and every $i \in \ZZ/e\ZZ$,
  \[
    \eps_i(\vke) \;=\; \max\{\,n \ge 0 : \ee_i^{\,n}\bigl(\vke \bmod q\Lcr\bigr) \ne 0\,\}
    \;\le\; 1.
  \]
\end{itemize}

The proof is elementary and Fock-native: it factors through a closed
expansion of $[\vke]$ in the crystal basis at $q = 0$, together with a
purely combinatorial signature-word analysis for partitions all of whose
parts are divisible by $e$. Gerber's bicrystal machinery~\cite{GN2017}
(which is stated only for $\ell \ge 2$) is not needed at level $1$; the
role played by his Cor.~5.3 in shaping the intuition is discussed in
Remark~\ref{rmk:gerber-shadow}. As a corollary we obtain that
$\vke \bmod q\Lcr$ is a $\ZZ$-linear combination of crystal-basis elements
$[e\lambda]$ indexed by partitions $\lambda \vdash k'$, with coefficients
equal to the number $f^\lambda$ of standard Young tableaux of shape $\lambda$
(Theorem~\ref{thm:expansion}).

We verify $\eps_i(\vke) \le 1$ empirically on $14$ target triples plus $15$
extended triples $(e,k',i)$, verify the expansion $[\vke] = \sum_\lambda
f^\lambda[e\lambda]$ on $17$ pairs $(e, k')$ (all with $e k' \le 16$), and
verify the individual bound $\eps_i(e\lambda) \le 1$ on $102$ partitions
$\lambda$ across ranges $e \in \{2,3,4,5\}$, $|\lambda| \le 7$ --- all pass.
\end{abstract}

\section{Setup}\label{sec:setup}

\subsection{Level-$1$ Uglov Fock space}

Fix $e \ge 2$. Let $\Fock_e$ denote the level-$1$ Uglov $q$-Fock space of
$\Uq$, viewed as a free $\ZZ[q, q^{-1}]$-module with standard basis
$\{|\lambda\rangle : \lambda \in \Pi\}$ indexed by all partitions
$\lambda = (\lambda_1 \ge \lambda_2 \ge \cdots \ge 0)$. The Chevalley
generators $e_i, f_i, K_i^{\pm 1}$ act on the standard basis by
Leclerc--Thibon's formulas (see~\cite{LT1996, Uglov2000}); we recall the
formulas for $e_i$ and $f_i$ that we need:
\begin{align*}
  f_i |\lambda\rangle
    &= \sum_{A \in \Aa_i(\lambda)} q^{N_i^r(\lambda, A)} \,|\lambda + A\rangle, \\
  e_i |\lambda\rangle
    &= \sum_{R \in \Rr_i(\lambda)} q^{-N_i^\ell(\lambda, R)} \,|\lambda - R\rangle,
\end{align*}
where $\Aa_i(\lambda)$ (resp.\ $\Rr_i(\lambda)$) is the set of addable
(resp.\ removable) $i$-nodes of $\lambda$, and $N_i^r$, $N_i^\ell$ are the
usual signature-based exponents. Residues of boxes are computed as
$\operatorname{res}(r, c) = c - r \bmod e$ using $0$-indexed rows and
columns.

\subsection{The ribbon operator $P_e$ and the vector $\vke$}

Define the \emph{ribbon-sign operator} $P_e$ on $\Fock_e$ by
\[
  P_e |\lambda\rangle \;=\; \sum_{\mu = \lambda + e\text{-ribbon}}
  (-q)^{h(\mu/\lambda)} |\mu\rangle,
\]
where the sum runs over all $\mu$ obtained from $\lambda$ by adding a
connected ribbon of size $e$, and $h(\mu/\lambda) = (\text{rows spanned}) - 1$
is the ribbon height. Set
\[
  \vke \;:=\; P_e^{k'} \,|\emptyset\rangle \;\in\; \Fock_e.
\]

\subsection{Kashiwara crystal lattice and operators}

Let $A := \ZZ_{(q)}$ denote the local ring at $q = 0$ of $\ZZ[q]$
(polynomials $g(q) \in \ZZ[q]$ localized by inverting those with
$g(0) \ne 0$).

\begin{definition}
The \emph{Kashiwara crystal lattice} of $\Fock_e$ is
\[
  \Lcr \;:=\; \bigoplus_{\lambda \in \Pi} A \cdot |\lambda\rangle
  \;\subset\; \Fock_e \otimes_{\ZZ[q, q^{-1}]} A[q^{-1}].
\]
The reduction $\Lcr/q\Lcr$ is a free $\ZZ$-module with basis
$\{[|\lambda\rangle] : \lambda \in \Pi\}$, called the crystal basis.
The Kashiwara operators $\ee_i, \ff_i$ act on $\Lcr/q\Lcr$ by the
good-node rule (Kleshchev--Uglov): for $\mu \in \Pi$,
\[
  \ee_i \,[\mu] \;=\; \begin{cases}
    [\mu \setminus \{n\}] & \text{if $\mu$ has a good removable $i$-node $n$,} \\
    0 & \text{otherwise,}
  \end{cases}
\]
where a good removable $i$-node is determined by reading the $i$-signature
of $\mu$ in bottom-up row order and cancelling adjacent $\Rr\Aa$ pairs; the
leftmost surviving $\Rr$ is the good one.\footnote{
We use $\Aa$ (Addable) and $\Rr$ (Removable) for signature symbols,
matching the code in \texttt{probe4\_canonical.py}.}
Extend $\ee_i, \ff_i$ to $\ZZ$-linear operators on $\Lcr/q\Lcr$.
\end{definition}

For $v \in \Lcr$, define
\[
  \eps_i(v) \;:=\; \max\{\, n \ge 0 : \ee_i^{\,n}\bigl([v]\bigr) \ne 0 \,\},
\]
with $\eps_i(0) := -\infty$.

That the good-node rule computes the Kashiwara operators on all of
$\Lcr/q\Lcr$ (not merely the sub-crystal $V(\Lam) \subset \Fock_e$) is
Ariki--Uglov; see~\cite{LT1996, Uglov2000}. In particular the operators
are $\ZZ$-linear, so
\begin{equation}\label{eq:linearity}
  \ee_i\Bigl(\sum_\mu c_\mu [\mu]\Bigr) \;=\; \sum_\mu c_\mu \,\ee_i[\mu].
\end{equation}

\section{Main theorem}

\begin{theorem}[C5 at level $1$]\label{thm:C5}
For every $e \ge 2$, every $k' \ge 1$, and every $i \in \ZZ/e\ZZ$,
\[
  \eps_i(\vke) \;\le\; 1.
\]
Concretely: $\ee_i^{\,2}\bigl([\vke]\bigr) = 0$ in $\Lcr/q\Lcr$.
\end{theorem}

The proof occupies Sections~\ref{sec:in-L}--\ref{sec:combine}.

\section{$\vke$ lies in the Kashiwara lattice}\label{sec:in-L}

\begin{proposition}\label{prop:in-L}
$\vke \in \Lcr$ for every $e \ge 2$ and $k' \ge 0$.
\end{proposition}

\begin{proof}
The Uglov ribbon formula
$P_e |\lambda\rangle = \sum_\mu (-q)^{h(\mu/\lambda)} |\mu\rangle$
has coefficients $(-q)^h$ with $h \ge 0$, so all coefficients lie in
$\ZZ[q] \subset A$. Since $|\emptyset\rangle \in \Lcr$ trivially and
$\Lcr$ is closed under $P_e$ (its matrix entries lie in $A$), $\vke$ is
obtained from $|\emptyset\rangle$ by $k'$ successive applications of $P_e$
and stays in $\Lcr$.
\end{proof}

\section{Closed form for $[\vke]$ modulo $q\Lcr$}\label{sec:expansion}

For a partition $\lambda = (\lambda_1, \dots, \lambda_\ell)$ let
$e\lambda := (e\lambda_1, e\lambda_2, \dots, e\lambda_\ell)$. Let
$f^\lambda$ denote the number of standard Young tableaux of shape $\lambda$
(equivalently, the number of chains
$\emptyset = \nu^{(0)} \subset \nu^{(1)} \subset \cdots \subset \nu^{(k')} = \lambda$
where each $\nu^{(t+1)}$ is obtained from $\nu^{(t)}$ by adding one box).

\begin{theorem}[Expansion of $\vke$ in the crystal basis at $q=0$]\label{thm:expansion}
For every $e \ge 2$ and $k' \ge 0$,
\[
  [\vke] \;=\; \sum_{\lambda \vdash k'} f^\lambda \,[\,|e\lambda\rangle\,]
  \qquad \text{in } \Lcr/q\Lcr.
\]
In particular, $\supp([\vke]) \subseteq \{e\lambda : \lambda \vdash k'\}$
inside the crystal basis.
\end{theorem}

\begin{proof}
Write $P_e^{(0)}$ for the specialization $P_e|_{q=0}$, acting on the
$\ZZ$-module $\Lcr/q\Lcr$. Since the Uglov formula weights each added
$e$-ribbon by $(-q)^h$, only the height-$0$ ribbons survive at $q = 0$:
\[
  P_e^{(0)} \,[|\lambda\rangle] \;=\; \sum_{\substack{\mu = \lambda + e\text{-ribbon} \\ h(\mu/\lambda) = 0}} [|\mu\rangle].
\]
A height-$0$ $e$-ribbon is a horizontal $e$-strip contained in a single row.

\emph{Restriction to partitions divisible by $e$.} We show by induction on $k'$
that the support of $P_e^{(0)^{k'}} \,[|\emptyset\rangle]$ lies in
$\{[|e\lambda\rangle] : \lambda \vdash k'\}$. Base $k' = 0$: support is
$\{[|\emptyset\rangle]\} = \{[|e \cdot \emptyset\rangle]\}$.
Inductive step: assume the support after $k'$ steps consists of $\mu = e\nu$
for various $\nu \vdash k'$. A height-$0$ $e$-ribbon added to $\mu = e\nu$
must be a single row of $e$ boxes; it can be added (a) by extending an
existing row $\mu_r$ from $e\nu_r$ to $e(\nu_r + 1)$, valid iff $r = 0$
or $\nu_{r-1} > \nu_r$; or (b) as a new row $(r = \ell(\nu), \text{length } e)$.
In case (a), the resulting partition has all parts divisible by $e$ and
equals $e\nu'$ where $\nu'$ is $\nu$ with $\nu_r$ incremented by $1$. In
case (b), the result equals $e\nu'$ where $\nu' = (\nu_1, \dots, \nu_{\ell(\nu)}, 1)$.
Either way the support after $k' + 1$ steps lies in
$\{e\nu' : \nu' \vdash k'+1\}$.

\emph{Coefficient of $[|e\lambda\rangle]$.} The multiplicity of $[|e\lambda\rangle]$
in $P_e^{(0)^{k'}} \,[|\emptyset\rangle]$ equals the number of sequences
$\nu^{(0)} = \emptyset,\, \nu^{(1)},\, \dots,\, \nu^{(k')} = \lambda$ of
partitions with $\nu^{(t+1)}$ obtained from $\nu^{(t)}$ by (i) extending
some existing row by one box in an admissible way (row $r = 0$ or
$\nu^{(t)}_{r-1} > \nu^{(t)}_r$), or (ii) appending a new row of length one.
But these are precisely the possible growth steps by an addable node of
$\nu^{(t)}$; the sequence is a standard Young tableau of shape $\lambda$
read in growth order. Hence the multiplicity is $f^\lambda$.
\end{proof}

\begin{example}\label{ex:expansions}
The small expansions:
\begin{align*}
  [\,v_{2,2}\,] &= [\,|(4)\rangle\,] + [\,|(2,2)\rangle\,] \\
  [\,v_{3,2}\,] &= [\,|(6)\rangle\,] + 2\,[\,|(4,2)\rangle\,] + [\,|(2,2,2)\rangle\,] \\
  [\,v_{2,3}\,] &= [\,|(6)\rangle\,] + [\,|(3,3)\rangle\,] \\
  [\,v_{3,3}\,] &= [\,|(9)\rangle\,] + 2\,[\,|(6,3)\rangle\,] + [\,|(3,3,3)\rangle\,] \\
  [\,v_{2,4}\,] &= [\,|(8)\rangle\,] + [\,|(4,4)\rangle\,]
\end{align*}
Coefficients are exactly $f^{(2)} = f^{(1,1)} = 1$, $f^{(3)} = f^{(1,1,1)} = 1$,
$f^{(2,1)} = 2$, matching Theorem~\ref{thm:expansion}. All five are computed
in \verb|probes/2026-08-11-C5-gerber/epsilon_i_probe.log|.
\end{example}

\section{The individual bound: $\eps_i([e\lambda]) \le 1$}\label{sec:individual}

\begin{proposition}[Individual crystal string bound]\label{prop:individual}
Let $\lambda$ be any partition and $i \in \ZZ/e\ZZ$. Then
$\eps_i\bigl([\,|e\lambda\rangle\,]\bigr) \le 1$; in fact
$\eps_i\bigl([\,|e\lambda\rangle\,]\bigr) \in \{0, 1\}$.
\end{proposition}

\begin{proof}
Set $\mu := e\lambda$, $L := \ell(\lambda)$. We analyze the addable and
removable $i$-nodes of $\mu$ and reduce the resulting $i$-signature
directly.

\medskip
\emph{Step 1: Row-residue classification.} With $0$-indexed rows and
columns, an addable node in row $r$ of $\mu$ (existing when the row is a
strict step, i.e., $r = 0$ or $\mu_{r-1} > \mu_r$) sits at column $\mu_r$
with residue $-r \bmod e$; a removable node in row $r$ (existing when
$r = L - 1$ or $\mu_{r+1} < \mu_r$) sits at column $\mu_r - 1$ with residue
$-r - 1 \bmod e$. In addition, if $r = L$ we may add a new row of length
$1$ with residue $-L \bmod e$. Since $e\lambda$ has all parts divisible by
$e$, the strict-step condition $\mu_{r-1} > \mu_r$ is equivalent to
$\lambda_{r-1} > \lambda_r$.

For fixed $i \in \ZZ/e\ZZ$:
\begin{itemize}
  \item candidate addable-$i$-rows: $r$ with $r \equiv -i \pmod e$ and $0 \le r \le L$;
  \item candidate removable-$i$-rows: $r$ with $r \equiv -(i+1) \pmod e$ and $0 \le r \le L - 1$.
\end{itemize}
Note that within any single row, addable and removable nodes carry
different residues (assuming $e \ge 2$), so each row contributes at most
one of $\Aa$, $\Rr$ to the $i$-signature.

\medskip
\emph{Step 2: Pairing structure.} The rows $r' \equiv -i - 1 \pmod e$ are
exactly the rows $r' = r - 1$ for $r$ a candidate addable-$i$-row with
$r \ge 1$. Thus every candidate addable-$i$-row $r \ge 1$ is immediately
preceded (in top-down order) by a candidate removable-$i$-row $r - 1$.
The residues match up in adjacent rows, and both the addable-at-$r$ and
removable-at-$r-1$ existence conditions reduce to the single strict-step
condition $\lambda_{r-1} > \lambda_r$.

We split into cases:

(a) \emph{$r = 0$, $i \equiv 0 \pmod e$.} The addable at $(0, \mu_0)$ is
always present (strict step at the top). No removable in ``row $-1$''
exists. This contributes a lone $\Aa$ symbol.

(b) \emph{$1 \le r \le L - 1$, $r \equiv -i \pmod e$.} Both addable-at-$r$
and removable-at-$r-1$ are present iff $\lambda_{r-1} > \lambda_r$. If
$\lambda_{r-1} = \lambda_r$, neither is present. So this row contributes
either $[\Aa, \Rr]$ (in top-down order at rows $r$ then $r - 1$? Actually
at rows $r$ (deeper) and $r - 1$ (shallower); in bottom-up order the
deeper row $r$ comes first) or nothing.

(c) \emph{$r = L$, $L \equiv -i \pmod e$.} The addable at the new row
$(L, 0)$ is present, and the removable-at-$r-1 = L - 1$ is present
(since $L - 1$ is the last row of $\mu$; the strict step at the bottom
$\mu_L = 0 < \mu_{L-1}$ is automatic). This contributes $[\Aa, \Rr]$ in
bottom-up order.

Between two consecutive candidate rows (which are separated by exactly
$e$ in the row index), rows of intermediate index carry residues in
$\ZZ/e\ZZ \setminus \{-i, -i-1\}$, contributing nothing to the
$i$-signature.

\medskip
\emph{Step 3: Bottom-up signature word.} Reading the $i$-signature of $\mu$
in bottom-up row order (largest row index first, matching
\verb|order='bottom_up'| in \verb|probe4_canonical.i_nodes|) gives a word
of the form
\[
  \Aa\Rr \,\cdot\, \Aa\Rr \,\cdots\, \Aa\Rr \,\cdot\, \Aa^?
\]
where each ``$\Aa\Rr$'' is contributed by case (b) or (c), and the
optional trailing ``$\Aa^?$'' is the lone $\Aa$ from case (a) (present
iff $i \equiv 0 \pmod e$).

\medskip
\emph{Step 4: Cancellation.} Kashiwara's cancellation rule for the
bottom-up signature is: while reading tokens left-to-right, cancel any
$\Aa$ that immediately follows an $\Rr$ on the stack. Concretely (in
pseudocode):
\begin{verbatim}
  for tok in seq:
      if tok == 'A' and stack and stack[-1] == 'R':
          stack.pop()      # cancel R-then-A
      else:
          stack.append(tok)
\end{verbatim}
The surviving stack always has $\Aa$'s at the bottom and $\Rr$'s on top
(no $\Rr$ can be immediately followed by an $\Aa$ that survives).
$\eps_i(\mu)$ equals the number of surviving $\Rr$'s.

Case A: the signature has $n \ge 0$ blocks $\Aa\Rr$ and \emph{no} trailing
$\Aa$. After processing the first block, the stack is $[\Aa, \Rr]$; each
subsequent $\Aa\Rr$ block first cancels the top $\Rr$ (leaving $[\Aa]$),
then pushes $\Rr$ (leaving $[\Aa, \Rr]$). By induction the final stack is
$[\Aa, \Rr]$ if $n \ge 1$, and empty if $n = 0$. Thus $\eps_i(\mu) \in \{0, 1\}$.

Case B: the signature has $n \ge 0$ blocks $\Aa\Rr$ followed by a trailing
$\Aa$. After processing the $\Aa\Rr$-blocks, the stack is $[\Aa, \Rr]$
(or empty if $n = 0$). Reading the trailing $\Aa$: it cancels the top $\Rr$,
leaving $[\Aa]$ (or if $n = 0$, pushing to $[\Aa]$). No surviving $\Rr$.
Thus $\eps_i(\mu) = 0$.

In both cases $\eps_i(\mu) \le 1$.
\end{proof}

\begin{remark}\label{rmk:zero-when-i-zero}
The proof shows $\eps_0([e\lambda]) = 0$ for every $\lambda$ (since $r = 0$
always contributes a trailing lone $\Aa$ when $i = 0$). Empirically this
matches all $17$ test cases (see
\verb|individual_and_expansion_check.log|).
\end{remark}

\begin{example}\label{ex:sig-2-4-2}
Take $e = 2$, $\lambda = (5, 3, 2, 1)$ so $\mu = e\lambda = (10, 6, 4, 2)$,
$L = 4$. Compute $\eps_1$:
\begin{itemize}
  \item candidate rows $r \equiv -1 \equiv 1 \pmod 2$: $r \in \{1, 3\}$;
  \item $r = 1$: $\lambda_0 > \lambda_1$ (5 > 3) $\Rightarrow$ block $[\Aa(1,6), \Rr(0,9)]$;
  \item $r = 3$: $\lambda_2 > \lambda_3$ (2 > 1) $\Rightarrow$ block $[\Aa(3,2), \Rr(2,3)]$.
\end{itemize}
Bottom-up signature: $[\Aa(3,2), \Rr(2,3), \Aa(1,6), \Rr(0,9)]$
$= \Aa\Rr\Aa\Rr$. Cancellation: push $\Aa$; push $\Rr$; $\Aa$ cancels
top $\Rr$; push $\Rr$. Final stack: $[\Aa, \Rr]$. So $\eps_1(\mu) = 1$,
with good removable node $(2, 3)$.
\end{example}

\section{Combining the pieces}\label{sec:combine}

\begin{proof}[Proof of Theorem~\ref{thm:C5}]
By Proposition~\ref{prop:in-L}, $\vke \in \Lcr$, so the reduction
$[\vke] \in \Lcr/q\Lcr$ is defined. By Theorem~\ref{thm:expansion},
\[
  [\vke] \;=\; \sum_{\lambda \vdash k'} f^\lambda \,[\,|e\lambda\rangle\,].
\]
By Proposition~\ref{prop:individual}, $\ee_i^{\,2}[\,|e\lambda\rangle\,] = 0$
for every $\lambda$. By $\ZZ$-linearity of $\ee_i^{\,2}$
(equation~\eqref{eq:linearity}),
\[
  \ee_i^{\,2}\bigl([\vke]\bigr)
  \;=\; \sum_{\lambda \vdash k'} f^\lambda \,\ee_i^{\,2}\bigl([\,|e\lambda\rangle\,]\bigr)
  \;=\; 0.
\]
Hence $\eps_i(\vke) \le 1$.
\end{proof}

\begin{remark}[Sharpness]\label{rmk:sharp}
The bound $\eps_i(\vke) \le 1$ is sharp for many $(e, k', i)$. For instance,
at $(e, k', i) = (2, 2, 1)$ we have $\ee_1([\vke]) = [|(3)\rangle] \ne 0$,
so $\eps_1(v_{2,2}) = 1$. See the table in Section~\ref{sec:verification}.
\end{remark}

\begin{remark}[Structural picture]\label{rmk:gerber-shadow}
Gerber--Norton~\cite{GN2017} prove a bicrystal commutation at level
$\ell \ge 2$ between the $\widehat{\mathfrak{sl}}_e$-Kashiwara operators and
their combinatorial $\mathfrak{sl}_\infty$-crystal (also called
the Heisenberg-crystal); their Corollary~5.3 predicts that in the
doubly-highest-weight component the elementary Heisenberg-crystal operators
send $|\emptyset, s\rangle \mapsto |\sigma[e], s\rangle$, i.e., produce
partitions with all parts divisible by $e$.
Their statements are formulated for $\ell \ge 2$; at $\ell = 1$ the
Heisenberg-crystal reduces to a single-runner abacus and the closed form
degenerates to $\tilde b_\sigma |\emptyset\rangle = |\sigma[e]\rangle$
componentwise. Theorem~\ref{thm:expansion} is exactly this closed form
promoted from crystal-basis elements to the class of the
\emph{$q$-quantum} vector $\vke = P_e^{k'}|\emptyset\rangle$: the $q$-Fock
computation at $q = 0$ reproduces the Heisenberg-crystal image of the
vacuum, weighted by SYT counts $f^\lambda$ that record the branching of
growth paths. In this sense the $\ell = 1$ case is the crystal shadow of
Gerber's Cor.~5.3; the extra $f^\lambda$ multiplicity is a genuinely level-$1$
phenomenon coming from the ordering of the $k'$ Heisenberg raisings that
Gerber packages into a single $\tilde b_\sigma$.

The bicrystal commutation, if applicable at level $1$ in a suitable
sense, would give $\eps_i(\vke) = 0$ (Heisenberg-descendant of a
Kashiwara-highest weight). What we prove is $\eps_i \le 1$, not
$\eps_i = 0$, and the difference is the ``controlled commutation
defect'' predicted at the quantum level by
Theorem~C4~\cite{Clio-C4-2026-08-14}: the identity
$[e_i, P_e] = (q - q^{-1})[e_i, C_e^{(1)}]$ vanishes at $q \to 0$ order
by order, but the $O(q^0)$ residue is non-trivial. The signature-word
proof of Proposition~\ref{prop:individual} shows this residue produces
at most one extra $\ee_i$-step per residue, exactly matching
$\eps_i \le 1$.
\end{remark}

\begin{remark}[Non-use of C4]\label{rmk:no-C4}
The proof of Theorem~\ref{thm:C5} is independent of
Theorem~C4~\cite{Clio-C4-2026-08-14}. C4 provides the quantum-level
$(q - q^{-1})$-divisibility of $e_i \cdot \vke$, and Path~2 of the PROVE.md
strategy proposed an inductive proof using $R_{i, k', e} := (e_i \vke)/(q-q^{-1})$.
That path is not needed: the direct combinatorial route through
Theorem~\ref{thm:expansion} and the signature analysis of
Proposition~\ref{prop:individual} bypasses the $R_{i,k',e}$ machinery.
An interesting side question, orthogonal to the current bound, is the
precise relation between $R_{i,k',e} \bmod q\Lcr$ and $\ee_i[\vke]$; we do
not pursue it here.
\end{remark}

\section{Computational verification}\label{sec:verification}

All computations are pure Python + integer arithmetic; no CAS required.
Code lives in \verb|~/projects/probes/2026-08-11-C5-gerber/| and reuses
\verb|qfock.py| (Laurent-polynomial arithmetic, standard basis,
$P_e |\lambda\rangle$) and \verb|probe4_canonical.py| (good-node
combinatorics for Kashiwara $\ee_i$, $\ff_i$).

\subsection{Phase 2 target set}

Compute $\eps_i(\vke)$ for the target set from
\verb|state/PROVE.md|:
\[
  (e, k') \in \{(2,2),\ (2,3),\ (3,2),\ (3,3),\ (4,2)\}, \qquad i \in \ZZ/e\ZZ.
\]

\begin{center}
\begin{tabular}{ccccc}
\toprule
$e$ & $k'$ & $i$ & $\eps_i(\vke)$ & one-step image \\
\midrule
2 & 2 & 0 & 0 & $-$ \\
2 & 2 & 1 & 1 & $|(3)\rangle$ \\
2 & 3 & 0 & 0 & $-$ \\
2 & 3 & 1 & 1 & $|(5)\rangle + 2|(3,2)\rangle + |(2,2,1)\rangle$ \\
3 & 2 & 0 & 0 & $-$ \\
3 & 2 & 1 & 1 & $|(3,2)\rangle$ \\
3 & 2 & 2 & 1 & $|(5)\rangle$ \\
3 & 3 & 0 & 0 & $-$ \\
3 & 3 & 1 & 1 & $2\,|(6,2)\rangle$ \\
3 & 3 & 2 & 1 & $|(8)\rangle + 2\,|(5,3)\rangle$ \\
4 & 2 & 0 & 0 & $-$ \\
4 & 2 & 1 & 0 & $-$ \\
4 & 2 & 2 & 1 & $|(4,3)\rangle$ \\
4 & 2 & 3 & 1 & $|(7)\rangle$ \\
\bottomrule
\end{tabular}
\end{center}

All $14$ triples satisfy $\eps_i \le 1$. Full log:
\verb|probes/2026-08-11-C5-gerber/epsilon_i_probe.log|.

\subsection{Extended $(e, k')$ sweep}

Same test on $(e, k') \in \{(2,4), (2,5), (3,4), (4,3), (5,2)\}$: all $15$
extra triples satisfy $\eps_i \le 1$. Max $\eps_i = 1$ throughout.
Log: \verb|probes/2026-08-11-C5-gerber/probe_C5_extended.log|.

\subsection{Verification of the expansion (Theorem~\ref{thm:expansion})}

For all $(e, k')$ with $e k' \le 16$ (a total of $17$ pairs across
$e \in \{2, 3, 4, 5\}$, $k' \in \{1, 2, 3, 4, 5\}$), verify that the class
$[\vke] \bmod q\Lcr$ equals $\sum_{\lambda \vdash k'} f^\lambda [e\lambda]$
by explicit computation (Python integer arithmetic).
All $17$ pairs pass. Log:
\verb|probes/2026-08-11-C5-gerber/individual_and_expansion_check.log|.

\subsection{Verification of the individual bound (Proposition~\ref{prop:individual})}

For all partitions $\lambda \vdash k'$ with $k' \le 7$ and $e \in \{2, 3, 4, 5\}$
(with $e k' \le 20$, giving $102$ distinct $\lambda$ across all $(e, k')$),
compute $\eps_i(e\lambda)$ for each $i \in \ZZ/e\ZZ$. All satisfy
$\eps_i \in \{0, 1\}$. Distribution: $191$ zero-values, $111$ one-values.
Log: \verb|probes/2026-08-11-C5-gerber/individual_and_expansion_check.log|.

\subsection{Signature-shape structural check}

Independent verification that the raw bottom-up $i$-signature of $e\lambda$
always has the form
$(\Aa\Rr)^n$ or $(\Aa\Rr)^n \Aa$
(equivalently: never has an $\Rr$ that is not immediately preceded, in
bottom-up order, by an $\Aa$). Sweep: $k' \in \{0, \dots, 7\}$,
$e \in \{2, 3, 4, 5\}$, all $\lambda \vdash k'$, all $i \in \ZZ/e\ZZ$.
$0$ failures. This directly certifies the proof structure of
Proposition~\ref{prop:individual}. Log:
\verb|probes/2026-08-11-C5-gerber/probe_C5_extended.log|.

\section{Consequences and open questions}

\subsection{§7 upgrade}

Combined with the previously-proved theorems T1, T2, T3, the $\kappa$-lemma
(2026-08-11), and C4 (2026-08-14), Theorem~\ref{thm:C5} promotes
Clio's §7 to \textbf{six proved theorems at level $1$}, with
\textbf{zero remaining open conjectures} at $\ell = 1$. Higher-level lifts
of C5 remain open (§\ref{sec:higher-level}).

\subsection{Higher-level lift}\label{sec:higher-level}

The proof of Proposition~\ref{prop:individual} is level-$1$-specific: it
uses that at level $1$ the abacus is a single runner, so ``rows of $\mu$
with all parts divisible by $e$'' has a clean row-by-row structure. At
level $\ell \ge 2$, one must work with $\ell$-abaci and
Gerber--Norton's Definition~3.5 (vessels, fore/aft periods) becomes
essential. We conjecture that a suitably-formulated
$\eps_i(\vke^{[\ell]}) \le 1$ statement holds at higher level, provable by
combining Gerber's arrow rule (Theorem~3.14 of~\cite{GN2017}) with a
higher-level analogue of Theorem~\ref{thm:expansion}. This is out of
scope for the current session.

\subsection{Refined form (part (b) of the PROVE.md statement)}

The PROVE.md refined form asked whether $R_{i, k', e} := (e_i \vke)/(q - q^{-1})$,
whose existence is guaranteed by C4, has class in $\Lcr/q\Lcr$ agreeing
with $\ee_i[\vke]$ up to a nonzero rational scalar. This is a natural
next question, not needed for the string-length bound, but likely provable
by combining C4's explicit $\Ce$ formula with the expansion of
Theorem~\ref{thm:expansion}. We flag it as a future PROVE candidate.

\begin{thebibliography}{9}

\bibitem{LT1996}
B.~Leclerc and J.-Y.~Thibon.
\newblock \emph{Canonical bases of $q$-deformed Fock spaces.}
\newblock Int.\ Math.\ Res.\ Not., 9:447--456, 1996.

\bibitem{Uglov2000}
D.~Uglov.
\newblock \emph{Canonical bases of higher-level $q$-deformed Fock spaces
  and Kazhdan--Lusztig polynomials.}
\newblock In: Physical combinatorics (Kyoto, 1999),
  Progr.\ Math., 191:249--299, Birkh\"auser, 2000.

\bibitem{GN2017}
T.~Gerber and E.~Norton.
\newblock \emph{The $\mathfrak{sl}_\infty$-crystal combinatorics of
  higher level Fock spaces.}
\newblock Transform.\ Groups 25 (2020), no.\ 4, 1119--1176. arXiv:1704.02169.

\bibitem{Iijima2012}
K.~Iijima.
\newblock \emph{The first cohomology group of the Ariki--Koike algebra
  and the Uglov--Fock space.}
\newblock arXiv:1207.6161, 2012.

\bibitem{Clio-C4-2026-08-14}
Clio (this project).
\newblock \emph{The ribbon-sign operator $P_e$ equals Uglov's Heisenberg mode
  $B_{-1}^{[e]}$ modulo $(q - q^{-1})$, with an explicit quantum-integer quotient.}
\newblock \verb|~/projects/proofs/2026-08-14-C4-iijima-B1.tex|, 2026.

\bibitem{Clio-commutation-defect-2026-08-12}
Clio (this project).
\newblock \emph{Heisenberg-commutant reframing of $v_{k',e}$ inside level-$1$
  Uglov Fock space.}
\newblock \verb|~/projects/memory/connections/2026-08-12-heisenberg-commutant-reframing.md|, 2026.

\end{thebibliography}

\end{document}
